% Pursuing Stacks (cotes 134-2 to 134-8): archive cote 134-4, site page 15 (PDF page 16), typescript page 103.
% typescript. Draft compared: ps3.tex:747-809.
% AI-transcribed by Claude Opus 5.5 (claude-opus-5-5, adaptive thinking, effort high), Message Batches API, batch ps1 (msgbatch_01Rv5hCaNTtUt4JHx1yqoHPL), 2026-11-12; not yet checked by eye.
\scanpage{134-4/15}
\mspage{103}
This I'll assume in the interpretation of $\underset{\underline{I}}{\hand{\sim}}$ in terms of strict homoto-%
pies. Namely, let
\[ \mathrm{Comp}(\underline{I}) \]
be the set of all homotopy intervals deduced from $\underline{I}$ by taking amalgama-%
tions of \over{\unsure{n}}{c}opies of $\underline{I}$ and $\overset{\hand{\vee}}{\underline{I}}$ , with an arbitrary number $n \hand{\geq} 1$ of summands.
Thus we get just $\underline{I}$ and $\overset{\hand{\vee}}{\underline{I}}$ for $n=1$, four intervals for \over{\unsure{m}}{n}$=2$, ..., $2^n$ inter-%
vals for $n$ arbitrary . It is now immediately checked that for $f,g$ in
$\mathrm{Hom}(X,Y)$, \xout{fxxxd} the relation $f \underset{\underline{I}}{\sim} g$ is equivalent to the existence of
$\underline{K} \in \mathrm{Comp}(\underline{I})$ , such that $f$ and $g$ be $\underline{K}$-homotopic (in the strict sense).

\uline{Remark} THe saturation conditions a')b')c') on \xout{xx} $W$ are easily checked
for the usual notion of weak equivalence for morphisms of topoi, and hence
also in the categories \xout{of the type} (Cat) and in any topos, and therefore
in any category $A^{\wedge}$\hand{$\hat{}$} (where it boils down \inserted{too} to the corresponding properties
on $W_{(\mathrm{Cat})}$ , as $W_A = i_A^{-1}(W_{(\mathrm{Cat})}$ ). Thus it \xout{\unsure{s}}seems definitely reasonabl
henceforth to take these as the standart notion of saturation (referring \xout{t}
to it's variants by the qualifications ``mild'' or ``strong''). On the other
hand, the stability conditions a)b)c) on $(M,W)$ are satisfied whenever $M$
is a topos, with the usual notion of weak equivalence - the condition
c') being a consequence of the more general \inserted[typed]{Mayer-Vietoris type} statement about amalgamations
of topoi under closed embeddings of such (cf. lemma on page ). The
same should hold in (Cat), with a similar Mayer-Vietoris argument -
there is a slight trouble here for applying the precedent result on
amalgamation of topoi, because a section $e \to C$ of an\xout{category} object $C$
of (Cat), namely an embedding of the one-point category $e=\hand{\Delta_0}$ into $C$ by
choice of an object of $C$, does not correspond in general to a closed
embedding of topoi. (In geometrical language, we get a ``point'' of the
topos $C^{\wedge}$ defined by $C$, but a point need not correspond to a subtopos,
let alone a closed one ...). This shows the asphericity criterion for
amalgamation of topoi, and hence also for amal\over{\unsure{ja}}{g}amation of categories,
has not been cut out with sufficient generality yet. As this whole
$\mathrm{Comp}(\underline{I})$ story is just a digression for the time being, I'll leave it at
that now.

M\over{\unsure{e}}{o}re important than amalgamation of intervals, is to compare the
notions of homotopy defined in terms of a\xout{n inter} homotopy interval $\underline{I}$, to
the corresponding notions for another interval\hand{,} $\underline{J}$ . Here the natural idea
first is to see what happens if we got a morphism of intervals (compatible
with endpoints, by definition)
\[ \underline{J} \to \underline{I} \quad . \]

\margin{15}
